\section{Contextualização e justificativa}
\label{sec-introducao-contextualizacao}

Pagamentos, compras parceladas e operações de crédito exigem a
interpretação de informações como saldo disponível, taxa de juros,
prazo e valor total a pagar. Ao investigar a alfabetização
financeira, Potrich, Vieira e Kirch consideram três dimensões:
conhecimento, atitude e comportamento financeiros
\cite{potrich2015}. No contexto escolar, uma situação de compra
permite articular o cálculo à decisão que ele informa. Saber
determinar o valor de uma prestação, por exemplo, é parte da
análise necessária para decidir se a compra cabe no orçamento.

Neste trabalho, apresenta-se o EDBank, um simulador bancário
desenvolvido para apoiar atividades de matemática financeira no
Ensino Médio. O participante utiliza uma conta fictícia vinculada
a uma agência, na qual pode receber um salário simulado, pagar
cobranças e realizar compras. Como essas operações afetam a mesma
conta, uma escolha feita durante a atividade modifica os recursos
disponíveis para a decisão seguinte.

O trabalho descreve como esse ambiente foi especificado e
implementado. O desenvolvimento compreende a definição das regras
das operações, a organização dos dados e a integração entre
interface e servidor. A finalidade educacional orienta a escolha
das situações representadas pelo simulador.

\section{Objetivo geral}
\label{sec-introducao-objetivo-geral}

Desenvolver uma aplicação web full stack que simule operações bancárias fictícias para apoiar o ensino de matemática financeira no Ensino Médio.


\section{Objetivos Específicos}
\label{sec-introducao-objetivo-especifico}

\begin{itemize}
    \item Especificar os requisitos funcionais e as regras de negócio das operações previstas para o simulador.

    \item Projetar a arquitetura da aplicação e modelar os dados necessários ao registro de contas e movimentações fictícias.

    \item Implementar os fluxos de abertura de conta, consulta de saldo e extrato, pagamentos, transferências e compras parceladas no ambiente simulado.

    \item Verificar o funcionamento das operações implementadas por meio de cenários de teste que contemplem cálculos financeiros, atualização de saldos e tratamento de operações inválidas.
\end{itemize}